\chapter{Dirichlet Characters and \texorpdfstring{$L$}{L}-Functions}

\begin{center}\itshape
In applying Fourier series to number theory, Dirichlet has recently found results that reach the summit of human acumen.\\
\upshape --- Carl Gustav Jacob Jacobi, letter to Alexander von Humboldt (1846)
\end{center}

Euler's proof that $\sum_p\frac1p$ diverges rests on the pole of $\zeta(s)=\prod_p(1-p^{-s})^{-1}$ at $s=1$. To count the primes $p\equiv a\pmod q$
separately, Dirichlet needed one product for each residue class, and there is none: the indicator of a class is not
multiplicative. His solution, in 1837, was to expand that indicator in the characters of the group $(\Z/q\Z)^*$ (Tome~2,
Proposition~11.15), each of which \emph{is} multiplicative and has an Euler product: the $L$-function $L(s,\chi)=\sum\chi(n)n^{-s}$.
The class of $a$ then gets its share of the divergence from the principal character, provided no other $L(s,\chi)$ kills the
pole by vanishing at $s=1$.

This chapter builds the $L$-functions and proves the decisive fact $L(1,\chi)\ne0$; the next chapter assembles Dirichlet's theorem. On
the way we meet primitive characters and Gauss sums, which control the size of character sums (Chapter~11), and an exact
formula for $L(1,\chi)$ that contains, in disguise, the class numbers of quadratic fields.

\begin{think}
For $\chi_4(n)=0,1,0,-1$ according as $n\equiv0,1,2,3\pmod4$, $L(1,\chi_4)=1-\frac13+\frac15-\cdots=\frac\pi4$ (Leibniz). Formally,
$\log L(s,\chi_4)\approx\sum_{p\equiv1(4)}p^{-s}-\sum_{p\equiv3(4)}p^{-s}$ as $s\to1^+$. What does $L(1,\chi_4)\ne0$ say about the primes in the two classes?
What would $L(1,\chi_4)=0$ have meant?
\end{think}

\section{Characters}

Fix $q\ge1$. A \emph{Dirichlet character modulo $q$} (Tome~2, Chapter~11) is a function $\chi:\Z\to\C$ with $\chi(n)=\psi(n\bmod q)$ when $\gcd(n,q)=1$, for a
group homomorphism $\psi:(\Z/q\Z)^*\to\C^*$, and $\chi(n)=0$ when $\gcd(n,q)>1$. It is $q$-periodic and completely multiplicative, its non-zero
values are $\varphi(q)$-th roots of unity, and there are $\varphi(q)$ characters modulo $q$, which form a group $\widehat G_q$ under multiplication; the
\emph{principal} character $\chi_0$ is the indicator of the integers prime to $q$. The inverse of $\chi$ is $\overline\chi$. A character is \emph{real} if
its values are real, that is $\chi^2=\chi_0$; otherwise it is \emph{complex}. It is \emph{even} or \emph{odd} according as $\chi(-1)=1$ or $-1$.

\begin{proposition}[orthogonality]\label{prop:orth}
\begin{enumerate}[label=(\alph*)]
\item For $\chi\ne\chi_0$, $\sum_{n=1}^q\chi(n)=0$, and $\bigl|\sum_{n\le N}\chi(n)\bigr|\le\varphi(q)$ for every $N\ge1$.
\item For $a$ prime to $q$ and every $n$, $\frac1{\varphi(q)}\sum_\chi\overline{\chi(a)}\chi(n)$ is $1$ if $n\equiv a\pmod q$, and $0$ otherwise.
\end{enumerate}
\end{proposition}
\begin{proof}
(a) Choose $b$ prime to $q$ with $\chi(b)\ne1$. As $n$ runs over a complete residue system, so does $bn$; hence
$\chi(b)\sum_n\chi(n)=\sum_n\chi(bn)=\sum_n\chi(n)$, and the sum vanishes. So a sum over $n\le N$ reduces to a sum over fewer than $q$
consecutive integers, at most $\varphi(q)$ of which are prime to $q$.
(b) is Tome~2, Proposition~11.15.
\end{proof}

\begin{example}\label{ex:smallchars}
Modulo $4$ there are two characters, $\chi_0$ and $\chi_4=\chi_{-4}$ above (odd). Modulo $3$: $\chi_0$ and $\chi_{-3}(n)=0,1,-1$ for $n\equiv0,1,2$ (odd). Modulo $5$
the group $(\Z/5\Z)^*$ is cyclic, generated by $2$, and a character is fixed by $\chi(2)\in\{1,i,-1,-i\}$; $\chi(2)=-1$ gives the Legendre symbol
$\chi_5(n)=\bigl(\frac n5\bigr)$, which is even, and $\chi(2)=\pm i$ gives two complex odd characters. For an odd prime $p$, the Legendre symbol
$\bigl(\frac\cdot p\bigr)$ is the unique character of order $2$ modulo $p$ (Tome~2, Chapter~11), and it is odd exactly when $p\equiv3\pmod4$, since
$\bigl(\frac{-1}p\bigr)=(-1)^{(p-1)/2}$.
\end{example}

Proposition~\ref{prop:orth}(a) is the property everything rests on: a non-principal character has bounded partial sums. They can still be
large compared with $\sqrt q$, and they are not random: Figure~\ref{fig:charsums} shows a bias that the end of the chapter explains.

\begin{figure}[ht]
\centering
\fig{ch03-charsums}
\caption{The partial sums $S(N)=\sum_{n\le N}\bigl(\frac n{1019}\bigr)$. Since $1019\equiv3\pmod4$, the symbol is odd and $S(1019-1-N)=S(N)$. The walk stays
almost entirely above $0$, and reaches $39=3\cdot13$ at $N=509$, where $13$ is the class number of $\Q(\sqrt{-1019})$
(Exercise~\ref{ex:residuesum} and the remark on class numbers).}
\label{fig:charsums}
\end{figure}

\section{Primitive characters}

The character $\chi_0\chi_{-4}$ modulo $12$ is $\chi_{-4}$ with the multiples of $3$ removed: it ``really'' has modulus $4$. To make this precise, for
$d\mid q$ let $U_d=\{n\in(\Z/q\Z)^*:n\equiv1\pmod d\}$, a subgroup.

\begin{lemma}\label{lem:lift}
Let $d\mid q$ and $m$ prime to $d$. There is an $n\equiv m\pmod d$ with $\gcd(n,q)=1$.
\end{lemma}
\begin{proof}
Let $t$ be the product of the primes dividing $q$ but not $m$, and $n=m+td$. A prime $p\mid q$ that divides $m$ does not divide $d$ (as
$\gcd(m,d)=1$) nor $t$, so $p\nmid td$ and $p\nmid n$. A prime $p\mid q$ that does not divide $m$ divides $t$, so again $p\nmid n$.
\end{proof}

\begin{proposition}\label{prop:conductor}
Let $\chi$ be a character modulo $q$, and let $D$ be the set of $d\mid q$ such that $\chi=1$ on $U_d$.
\begin{enumerate}[label=(\alph*)]
\item If $d_1,d_2\in D$ then $\gcd(d_1,d_2)\in D$. So the least element $d^*$ of $D$, the \emph{conductor} of $\chi$, divides every element of $D$.
\item For $d\in D$ there is a unique character $\psi$ modulo $d$ with $\chi(n)=\psi(n)$ for all $n$ prime to $q$; we say $\psi$ \emph{induces} $\chi$. Conversely,
if a character modulo $d\mid q$ induces $\chi$, then $d\in D$.
\item Call $\chi$ \emph{primitive} if $d^*=q$. Exactly one primitive character induces $\chi$: the character $\chi^*$ modulo $d^*$ given by (b).
\end{enumerate}
\end{proposition}
\begin{proof}
(a) Write $q=\prod p^{e_p}$. By the Chinese remainder theorem $(\Z/q\Z)^*\cong\prod_p(\Z/p^{e_p}\Z)^*$, and $n\equiv1\pmod d$ means $n\equiv1\pmod{p^{v_p(d)}}$
for every $p$. So $U_d=\prod_pV_p(v_p(d))$, where $V_p(f)$ is the subgroup of residues $\equiv1\pmod{p^f}$ in $(\Z/p^{e_p}\Z)^*$; these subgroups decrease
as $f$ grows. Hence $U_{d_1}U_{d_2}=\prod_pV_p(\min(v_p(d_1),v_p(d_2)))=U_{\gcd(d_1,d_2)}$, and $\chi$, trivial on $U_{d_1}$ and $U_{d_2}$, is trivial on their product.
If $d\in D$, then $\gcd(d,d^*)\in D$ and $\gcd(d,d^*)\le d^*$, so $\gcd(d,d^*)=d^*$, that is $d^*\mid d$.

(b) For $m$ prime to $d$, pick $n$ as in Lemma~\ref{lem:lift} and set $\psi(m)=\chi(n)$. If $n'$ is another choice, $n'\equiv n\pmod d$, so $n'n^{-1}\in U_d$ (inverse
modulo $q$) and $\chi(n')=\chi(n)$: $\psi$ is well defined, and it is multiplicative because $nn'$ is a valid choice for $mm'$. Extended by $0$ off the
residues prime to $d$, $\psi$ is a character modulo $d$ that agrees with $\chi$ on integers prime to $q$; it is unique because, by the lemma, every
residue prime to $d$ contains such integers. Conversely, if $\psi$ modulo $d$ induces $\chi$, then for $n\in U_d$, $\chi(n)=\psi(n)=\psi(1)=1$.

(c) Let $\chi^*$ be the character modulo $d^*$ inducing $\chi$. If $\chi^*$ were trivial on its own subgroup $U'_{d'}$ for some $d'\mid d^*$, $d'<d^*$,
then $\chi$ would be trivial on $U_{d'}$ (for $n\in U_{d'}$, $\chi(n)=\chi^*(n)=1$), contradicting the minimality of $d^*$; so $\chi^*$ is primitive. Let $\psi$
modulo $d$ be primitive and induce $\chi$. By (b), $d\in D$, so $d^*\mid d$ by (a). For $m$ prime to $d$ with $m\equiv1\pmod{d^*}$, lift $m$ to $n$ prime to $q$
(Lemma~\ref{lem:lift}); then $n\in U_{d^*}$ and $\psi(m)=\chi(n)=1$. So $\psi$ is trivial on the subgroup attached to $d^*$ modulo $d$, and primitivity
forces $d=d^*$; then $\psi=\chi^*$ by the uniqueness in (b).
\end{proof}

For example, $\chi_{-4}$, $\chi_{-3}$ and $\chi_5$ are primitive, and the conductor of $\chi_0$ is $1$. For a prime modulus every non-principal
character is primitive, since the only proper divisor is $1$ and $U_1$ is the whole group.

\section{Gauss sums}

Write $e(x)=e^{2\pi ix}$. The \emph{Gauss sum} of a character $\chi$ modulo $q$ is
\[
\tau(\chi)=\sum_{a=1}^q\chi(a)e\Bigl(\frac aq\Bigr).
\]
For the Legendre symbol modulo $p$ it is the quadratic Gauss sum of Tome~4, Proposition~7.13, whose value Gauss determined
(Tome~4, Theorem~8.14).

\begin{theorem}\label{thm:gauss}
Let $\chi$ be a primitive character modulo $q$. Then for every integer $n$,
\[
\sum_{a=1}^q\chi(a)e\Bigl(\frac{an}q\Bigr)=\overline{\chi(n)}\,\tau(\chi),
\]
and $|\tau(\chi)|=\sqrt q$.
\end{theorem}
\begin{proof}
If $\gcd(n,q)=1$, substitute $a=n^{-1}b$ (inverse modulo $q$): as $a$ runs over the residues, so does $b$, and
$\chi(a)=\overline{\chi(n)}\chi(b)$ since $\chi(n)\chi(n^{-1})=1$ and $|\chi(n)|=1$. The sum becomes $\overline{\chi(n)}\sum_b\chi(b)e(\frac bq)$.

If $g=\gcd(n,q)>1$, the right side is $0$, and we show that the left side $T$ vanishes too. Write $q=gq_1$, $n=gn_1$, so $e(\frac{an}q)=e(\frac{an_1}{q_1})$
depends only on $a\bmod q_1$. Since $q_1<q$ divides $q$ and $\chi$ is primitive, $q_1\notin D$: there is $c\in U_{q_1}$ with $\chi(c)\ne1$. As $a$ runs over the residues
modulo $q$, so does $ca$, and $ca\equiv a\pmod{q_1}$; hence $T=\sum_a\chi(ca)e(\frac{an_1}{q_1})=\chi(c)T$, and $T=0$.

For the modulus, let $T$ be the sum over $n=1,\dots,q$ of the squared moduli of both sides. On the right we get $T=\varphi(q)|\tau(\chi)|^2$. On the left, expanding,
\[
\begin{aligned}T&=\sum_{a,b}\chi(a)\overline{\chi(b)}\sum_{n=1}^qe\Bigl(\frac{(a-b)n}q\Bigr)\eqbr=q\sum_a|\chi(a)|^2=q\varphi(q),\end{aligned}
\]
because the inner sum is $q$ if $a=b$ and $0$ otherwise (a geometric sum of a $q$-th root of unity $\ne1$). So $|\tau(\chi)|^2=q$.
\end{proof}

So a primitive character is, up to the constant $\tau(\chi)$, its own Fourier transform on $\Z/q\Z$: its spectrum is flat. This is the
character analogue of Tome~4, Proposition~7.13, and the key to the Pólya--Vinogradov inequality of Chapter~11. Here it gives a formula for $L(1,\chi)$.

\section{\texorpdfstring{$L$}{L}-functions}

The \emph{Dirichlet $L$-function} of $\chi$ is $L(s,\chi)=\sum_{n\ge1}\chi(n)n^{-s}$. We write $s=\sigma+it$.

\begin{theorem}\label{thm:lbasic}
Let $\chi$ be a character modulo $q$.
\begin{enumerate}[label=(\alph*)]
\item For $\sigma>1$ the series converges absolutely, $L(s,\chi)=\prod_p\bigl(1-\chi(p)p^{-s}\bigr)^{-1}$, and
\[
L(s,\chi)=\exp\Bigl(\sum_p\sum_{k\ge1}\frac{\chi(p)^k}{kp^{ks}}\Bigr).
\]
In particular $L(s,\chi)\ne0$ for $\sigma>1$.
\item If $\chi\ne\chi_0$, the series converges on $\sigma>0$ and $L(s,\chi)$ is holomorphic there.
\item $L(s,\chi_0)=\zeta(s)\prod_{p\mid q}(1-p^{-s})$ for $\sigma>1$, so $L(s,\chi_0)-\frac{\varphi(q)/q}{s-1}$ extends holomorphically to $\sigma>0$.
\end{enumerate}
\end{theorem}
\begin{proof}
(a) $|\chi(n)n^{-s}|\le n^{-\sigma}$, and $\chi$ is completely multiplicative, so Tome~4, Theorem~10.3 gives the Euler product, with factors
$\sum_k\chi(p)^kp^{-ks}=(1-\chi(p)p^{-s})^{-1}$. For $|w|<1$, $\exp\bigl(\sum_{k\ge1}\frac{w^k}k\bigr)=\frac1{1-w}$: both sides are holomorphic on the unit disc, and the
left side $E(w)$ satisfies $E'=\frac E{1-w}$, $E(0)=1$, so $((1-w)E)'=0$. With $w=\chi(p)p^{-s}$, the partial products of the Euler product over $p\le y$
are $\exp\bigl(\sum_{p\le y}\sum_k\frac{\chi(p)^k}{kp^{ks}}\bigr)$. The double series converges absolutely (its terms are bounded by those of
$\sum_p\sum_kp^{-k\sigma}<\infty$), so letting $y\to\infty$ and using the continuity of $\exp$ gives the formula. An exponential never vanishes.

(b) The partial sums of $\chi$ are bounded by Proposition~\ref{prop:orth}(a); apply Tome~4, Proposition~10.1(b).

(c) $\chi_0(n)=1$ exactly when $n$ has no prime factor dividing $q$, so the Euler product of $L(s,\chi_0)$ is that of $\zeta$ without the factors at
$p\mid q$. By Tome~4, Proposition~10.6, $\zeta(s)-\frac1{s-1}$ is holomorphic on $\sigma>0$, and $P(s)=\prod_{p\mid q}(1-p^{-s})$ is entire, so
$L(s,\chi_0)-\frac{P(1)}{s-1}=P(s)\bigl(\zeta(s)-\frac1{s-1}\bigr)+\frac{P(s)-P(1)}{s-1}$ extends holomorphically, and $P(1)=\prod_{p\mid q}(1-\frac1p)=\frac{\varphi(q)}q$.
\end{proof}

If $\chi^*$ modulo $d$ induces $\chi$ modulo $q$, the Euler products differ only at the primes dividing $q$:
$L(s,\chi)=L(s,\chi^*)\prod_{p\mid q}(1-\chi^*(p)p^{-s})$. So questions about $L(s,\chi)$ near $s=1$ reduce to primitive characters, up to finitely many
explicit factors that do not vanish at $s=1$.

\section{The non-vanishing of \texorpdfstring{$L(1,\chi)$}{L(1,chi)}}

Everything now hinges on one fact. Figure~\ref{fig:lreal} shows three examples.

\begin{figure}[ht]
\centering
\fig{ch03-lreal}
\caption{$L(\sigma,\chi)$ for $0<\sigma\le3$ and the real characters $\chi_{-3}$, $\chi_{-4}$, $\chi_5$. At $\sigma=1$ the values are $\frac\pi{3\sqrt3}=0.6046$, $\frac\pi4=0.7854$ and
$\frac2{\sqrt5}\log\frac{1+\sqrt5}2=0.4304$ (Theorem~\ref{thm:l1formula}). As $\sigma\to\infty$ all tend to $\chi(1)=1$. The even character $\chi_5$ tends to $0$ as $\sigma\to0$.}
\label{fig:lreal}
\end{figure}

\begin{lemma}\label{lem:prodchars}
For real $\sigma>1$, $\prod_{\chi\bmod q}L(\sigma,\chi)\ge1$.
\end{lemma}
\begin{proof}
By Theorem~\ref{thm:lbasic}(a), the product is $\exp(S)$ with
\[
\begin{aligned}S&=\sum_p\sum_{k\ge1}\frac1{kp^{k\sigma}}\sum_\chi\chi(p^k)\eqbr=\varphi(q)\sum_{p^k\equiv1\,(q)}\frac1{kp^{k\sigma}}\ge0,\end{aligned}
\]
using Proposition~\ref{prop:orth}(b) with $a=1$ (the sums may be exchanged, the double series converging absolutely). So $S$ is real and
non-negative, and $\exp(S)\ge1$.
\end{proof}

\begin{theorem}[Dirichlet, 1837]\label{thm:lnonzero}
$L(1,\chi)\ne0$ for every character $\chi\ne\chi_0$.
\end{theorem}
\begin{proof}
\emph{Complex characters.} Suppose $L(1,\chi)=0$ with $\chi\ne\overline\chi$. Conjugating the convergent series, $L(1,\overline\chi)=\overline{L(1,\chi)}=0$, and $\overline\chi$ is
another factor of the product. As $\sigma\to1^+$: $0<L(\sigma,\chi_0)\le\zeta(\sigma)\le\frac\sigma{\sigma-1}$ (comparing $\sum n^{-\sigma}$ with $1+\int_1^\infty t^{-\sigma}\dd t$); each of
$L(\sigma,\chi)$ and $L(\sigma,\overline\chi)$ is $O(\sigma-1)$, being holomorphic near $1$ with a zero there; and the other factors are bounded, being continuous
at $1$. So the product is $O\bigl(\frac1{\sigma-1}(\sigma-1)^2\bigr)\to0$, contradicting Lemma~\ref{lem:prodchars}.

\emph{Real characters.} Suppose $L(1,\chi)=0$ with $\chi$ real, $\chi\ne\chi_0$. Let $r=\mathbf1*\chi$, so $r(n)=\sum_{d\mid n}\chi(d)$. It is multiplicative, and at a
prime power $r(p^k)=\sum_{j=0}^k\chi(p)^j$ equals $k+1$ if $\chi(p)=1$, $1$ if $\chi(p)=0$, and $1$ or $0$ according as $k$ is even or odd if $\chi(p)=-1$.
Hence $r(n)\ge0$ for all $n$, and $r(m^2)\ge1$. By Tome~4, Proposition~10.2(a), $F(s)=\zeta(s)L(s,\chi)=\sum r(n)n^{-s}$ for $\sigma>1$.

Now $\zeta(s)-\frac1{s-1}$ is holomorphic on $\sigma>0$, and so is $\frac{L(s,\chi)}{s-1}$ (extended by $L'(1,\chi)$ at $s=1$), since $L(1,\chi)=0$. So
$F(s)=\bigl(\zeta(s)-\frac1{s-1}\bigr)L(s,\chi)+\frac{L(s,\chi)}{s-1}$ extends holomorphically to $\sigma>0$, which contains the disc $|s-2|<2$. By Tome~4,
Theorem~9.11, $F$ is the sum of its Taylor series at $2$ on that disc. By Tome~4, Proposition~10.1(a) applied $k$ times, $F^{(k)}(2)=\sum_nr(n)(-\log n)^kn^{-2}$. At $s=\frac12$,
\[
\begin{aligned}F\Bigl(\frac12\Bigr)&=\sum_{k\ge0}\frac{F^{(k)}(2)}{k!}\Bigl(-\frac32\Bigr)^k\eqbr=\sum_{k\ge0}\sum_{n\ge1}\frac{r(n)}{n^2}\frac{(\frac32\log n)^k}{k!}\eqbr=\sum_{n\ge1}\frac{r(n)}{n^2}e^{\frac32\log n}\eqbr=\sum_{n\ge1}\frac{r(n)}{\sqrt n},\end{aligned}
\]
where the exchange of the two sums is allowed because all terms are non-negative. So $\sum r(n)n^{-1/2}$ converges. But it is at least
$\sum_mr(m^2)m^{-1}\ge\sum_m\frac1m=\infty$. This contradiction proves the theorem.
\end{proof}

\begin{beauty}
The two halves of the proof use positivity in two different ways. For complex $\chi$, the product of all $L$-functions is the
generating function of the primes $p\equiv1\pmod q$, whose coefficients are non-negative: two zeros would beat one pole. For real $\chi$, only one
zero is available, and the positive object is $\zeta(s)L(s,\chi)$, the zeta function of the quadratic field attached to $\chi$; its coefficients
$r(n)$ count ideals of norm $n$. A Dirichlet series with non-negative coefficients cannot be holomorphic beyond the point where it
diverges (Landau's theorem, which the proof reproves in this special case), and the squares make it diverge at $\frac12$. The real case is the
deep one: the question of how small $L(1,\chi)$ can be for real $\chi$, the problem of \emph{Siegel zeros}, is still open.
\end{beauty}

\section{The value of \texorpdfstring{$L(1,\chi)$}{L(1,chi)}}

Dirichlet did better than $L(1,\chi)\ne0$: he computed $L(1,\chi)$ in finite terms. We need the sum of the series $\sum\frac{z^n}n$ on the unit circle.

\begin{lemma}\label{lem:logseries}
For $0<x<2\pi$, $\displaystyle\sum_{n\ge1}\frac{e^{inx}}n=-\log\Bigl(2\sin\frac x2\Bigr)+i\,\frac{\pi-x}2$.
\end{lemma}
\begin{proof}
Let $z=e^{ix}\ne1$. For $t\in[0,1]$, $1-zt\ne0$ (as $|zt|=1$ forces $t=1$, and $z\ne1$), so $|1-zt|\ge c>0$ on $[0,1]$ by compactness. Then
\[
\sum_{n=1}^N\frac{z^n}n=\int_0^1\sum_{n=1}^Nz^nt^{n-1}\dd t=\int_0^1\frac{z-z^{N+1}t^N}{1-zt}\dd t,
\]
and the part with $t^N$ has modulus at most $\frac1c\int_0^1t^N\dd t\to0$. So the series converges to $\int_0^1\frac z{1-zt}\dd t$. Now $\operatorname{Re}(1-zt)=1-t\cos x>0$
for $t\in[0,1[$, so $-\operatorname{Log}(1-zt)$ (principal logarithm) is a $C^1$ primitive of $\frac z{1-zt}$ on $[0,1]$, and the integral is $-\operatorname{Log}(1-z)$. Finally
$1-e^{ix}=-2i\sin\frac x2\,e^{ix/2}=2\sin\frac x2\,e^{i(x-\pi)/2}$, with $2\sin\frac x2>0$ and $\frac{x-\pi}2\in\left]-\frac\pi2,\frac\pi2\right[$, so $\operatorname{Log}(1-z)=\log(2\sin\frac x2)+i\frac{x-\pi}2$.
\end{proof}

Taking imaginary parts recovers $\sum\frac{\sin nx}n=\frac{\pi-x}2$ (Tome~4, Chapter~6); the real part gives $\sum\frac{\cos nx}n=-\log(2\sin\frac x2)$.

\begin{theorem}[Dirichlet's formula]\label{thm:l1formula}
Let $\chi$ be a primitive character modulo $q>1$. Then
\[
\begin{gathered}L(1,\chi)=-\frac{i\pi}{q\,\tau(\overline\chi)}\sum_{a=1}^{q-1}\overline\chi(a)\,a\quad(\chi\text{ odd}),\\
L(1,\chi)=-\frac1{\tau(\overline\chi)}\sum_{a=1}^{q-1}\overline\chi(a)\log\sin\frac{\pi a}q\quad(\chi\text{ even}).\end{gathered}
\]
\end{theorem}
\begin{proof}
$\overline\chi$ is primitive with $|\tau(\overline\chi)|=\sqrt q$, so Theorem~\ref{thm:gauss} for $\overline\chi$ gives $\chi(n)=\frac1{\tau(\overline\chi)}\sum_{a=1}^{q-1}\overline\chi(a)e(\frac{an}q)$ (the term $a=q$ vanishes,
as $\overline\chi(q)=0$ for $q>1$). Insert this into $L(1,\chi)=\sum_n\frac{\chi(n)}n$ and exchange the finite sum over $a$ with the convergent series over $n$. By
Lemma~\ref{lem:logseries} with $x=\frac{2\pi a}q$,
\[
\begin{aligned}L(1,\chi)&=\frac1{\tau(\overline\chi)}\sum_{a=1}^{q-1}\overline\chi(a)\Bigl(-\log\Bigl(2\sin\frac{\pi a}q\Bigr)\eqbr+i\pi\Bigl(\frac12-\frac aq\Bigr)\Bigr).\end{aligned}
\]
Since $\sum_a\overline\chi(a)=0$, the constants $\log2$ and $\frac{i\pi}2$ drop out. If $\chi$ is odd, so is $\overline\chi$, and the substitution $a\mapsto q-a$ shows
$\sum_a\overline\chi(a)\log\sin\frac{\pi a}q=-\sum_a\overline\chi(a)\log\sin\frac{\pi a}q=0$, leaving the first formula. If $\chi$ is even, the same substitution shows
$\sum_a\overline\chi(a)a=\sum_a\overline\chi(a)(q-a)=-\sum_a\overline\chi(a)a$, so this sum vanishes, leaving the second.
\end{proof}

\begin{example}\label{ex:leibniz}
For $\chi_{-4}$: $\overline\chi=\chi$, $\tau(\chi_{-4})=e(\frac14)-e(\frac34)=2i$, and $\sum_a\chi(a)a=1-3=-2$, so $L(1,\chi_{-4})=-\frac{i\pi}{4\cdot2i}(-2)=\frac\pi4$, Leibniz's series.
\end{example}

For a real primitive character, $\tau(\chi)^2=\chi(-1)q$ (Exercise~\ref{ex:gausssquare}), so $\tau(\chi)=\pm\sqrt q$ or $\pm i\sqrt q$; Gauss's theorem fixes the sign
for the Legendre symbol: $\tau=\sqrt p$ if $p\equiv1\pmod4$ and $i\sqrt p$ if $p\equiv3\pmod4$ (Tome~4, Theorem~8.14). For $p\equiv3\pmod4$ the first formula
becomes $L(1,\chi)=-\frac\pi{p^{3/2}}\sum_{a<p}\bigl(\frac ap\bigr)a$, and since $L(1,\chi)>0$ (Exercise~\ref{ex:lpositive}) we get a purely arithmetic statement with no
known elementary proof (Exercise~\ref{ex:residuesum}).

\begin{remark}[for culture]
For $p\equiv3\pmod4$, $p>3$, the number $h(-p)=-\frac1p\sum_{a<p}\bigl(\frac ap\bigr)a$ is a positive integer, the \emph{class number} of $\Q(\sqrt{-p})$: it
measures the failure of unique factorisation in its ring of integers (Tome~1, Chapter~11; Tome~12). Then $L(1,\chi)=\frac{\pi h(-p)}{\sqrt p}$, Dirichlet's
\emph{class number formula}, and $\sum_{n<p/2}\bigl(\frac np\bigr)=\bigl(2-\bigl(\frac2p\bigr)\bigr)h(-p)$, which is the maximum of the walk in Figure~\ref{fig:charsums}. For
even characters the second formula produces logarithms of units of real quadratic fields, as $\log\frac{1+\sqrt5}2$ for $q=5$. Proofs are in
Davenport, Chapter~6, and in Tome~12.
\end{remark}

\section{Exercises}

\begin{exercise}[$\star$]\label{ex:chartables}
Write the tables of the characters modulo $5$, $8$ and $12$. Which are real, which are odd, and which are primitive? Give the conductor of
each.
\end{exercise}

\begin{exercise}[$\star\star$]\label{ex:primcount}
Let $\varphi^*(q)$ be the number of primitive characters modulo $q$. Show that $\sum_{d\mid q}\varphi^*(d)=\varphi(q)$, deduce that $\varphi^*=\mu*\varphi$ is multiplicative
with $\varphi^*(p)=p-2$ and $\varphi^*(p^k)=p^k(1-\frac1p)^2$ for $k\ge2$, and that there are no primitive characters modulo $q\equiv2\pmod4$.
\end{exercise}

\begin{exercise}[$\star\star$]\label{ex:gausssquare}
Show that $\overline{\tau(\chi)}=\chi(-1)\tau(\overline\chi)$ for every character, and that $\tau(\chi)\tau(\overline\chi)=\chi(-1)q$ when $\chi$ is primitive. Deduce that
$\tau(\chi)^2=\chi(-1)q$ for a real primitive character.
\end{exercise}

\begin{exercise}[$\star$]\label{ex:lpositive}
Show that $L(1,\chi)>0$ for every real character $\chi\ne\chi_0$.
\end{exercise}

\begin{exercise}[$\star\star$]\label{ex:l1values}
\emph{Signature exercise.} Show that $L(1,\chi_{-3})=1-\frac12+\frac14-\frac15+\frac17-\cdots=\frac\pi{3\sqrt3}$ and
$L(1,\chi_5)=1-\frac12-\frac13+\frac14+\frac16-\cdots=\frac2{\sqrt5}\log\frac{1+\sqrt5}2$.
\end{exercise}

\begin{exercise}[$\star\star$]\label{ex:residuesum}
Let $p\equiv3\pmod4$ be a prime. Show that the sum of the quadratic non-residues in $\{1,\dots,p-1\}$ exceeds the sum of the quadratic
residues. Check it for $p=7$ and $p=11$.
\end{exercise}

\begin{exercise}[$\star\star$]\label{ex:lderivative}
Let $q\ge3$ and $\chi\ne\chi_0$ modulo $q$. Show that $L'(s,\chi)=-\sum_n\chi(n)(\log n)n^{-s}$ on $\sigma>0$, that $|L(1,\chi)|\le\log q+2$, and that
$|L'(1,\chi)|\le\frac12\log^2q+\log q+1$. (Split the series at $n=q$ and use Abel summation on the tail.)
\end{exercise}

\section{Hints}

\begin{hint}{ex:chartables}
$(\Z/8\Z)^*\cong(\Z/2\Z)^2$ and $(\Z/12\Z)^*\cong(\Z/2\Z)^2$; a character is fixed by its values at generators. Use Proposition~\ref{prop:conductor} to test
the proper divisors $6$, $4$, $3$, $2$, $1$.
\end{hint}
\begin{hint}{ex:primcount}
Each character modulo $q$ is induced by exactly one primitive character modulo some $d\mid q$, and conversely each primitive character
modulo $d\mid q$ induces exactly one character modulo $q$. Then use Möbius inversion.
\end{hint}
\begin{hint}{ex:gausssquare}
Conjugate the definition and substitute $a\mapsto-a$. Then combine with $|\tau(\chi)|^2=q$.
\end{hint}
\begin{hint}{ex:lpositive}
$L(\sigma,\chi)>0$ for real $\sigma>1$, and $L(\sigma,\chi)$ is continuous on $\left]0,\infty\right[$.
\end{hint}
\begin{hint}{ex:l1values}
Theorem~\ref{thm:l1formula}: compute $\tau(\chi_{-3})$ and $\tau(\chi_5)$ directly, and $\sin\frac{2\pi}5=2\sin\frac\pi5\cos\frac\pi5$ with $2\cos\frac\pi5=\frac{1+\sqrt5}2$.
\end{hint}
\begin{hint}{ex:residuesum}
Combine the first formula of Theorem~\ref{thm:l1formula} with Gauss's sign and Exercise~\ref{ex:lpositive}.
\end{hint}
\begin{hint}{ex:lderivative}
For $n\le q$ bound each term trivially. For $n>q$, write the tail with $A(t)=\sum_{n\le t}\chi(n)$, $|A(t)|\le\varphi(q)\le q$.
\end{hint}

\section{Solutions}

\begin{solution}{ex:chartables}
\idea A character of a product of cyclic groups is a choice of a root of unity of the right order at each generator. The conductor is the
least $d$ such that the character is $1$ on all units $\equiv1\pmod d$.
\writeup Modulo $5$: generator $2$, $\chi(2)\in\{1,i,-1,-i\}$, giving $\chi_0$, two complex odd characters ($\chi(4)=\chi(2)^2=-1$, and $4\equiv-1$), and $\chi_5$, real
and even. All three non-principal ones are primitive ($5$ is prime). Modulo $8$: $(\Z/8\Z)^*=\{1,3,5,7\}$, every element of order $2$, generated by
$3$ and $5$; the four characters are real: $\chi_0$; $\chi_{-4}$ ($1,-1,1,-1$ on $1,3,5,7$; conductor $4$, odd); $\chi_8$ ($1,-1,-1,1$: $\chi_8(n)=1$ iff $n\equiv\pm1\pmod8$;
conductor $8$, even); $\chi_{-8}=\chi_{-4}\chi_8$ ($1,1,-1,-1$; conductor $8$, odd). Modulo $12$: $(\Z/12\Z)^*=\{1,5,7,11\}$; the characters are $\chi_0$, the one induced by
$\chi_{-4}$ ($1,1,-1,-1$; conductor $4$, odd), the one induced by $\chi_{-3}$ ($1,-1,1,-1$; conductor $3$, odd), and their product $\chi_{12}$ ($1,-1,-1,1$), which is
not trivial on $U_6=U_3=\{1,7\}$, $U_4=\{1,5\}$, $U_2$ or $U_1$: it is primitive, even, with $\chi_{12}(n)=1$ iff $n\equiv\pm1\pmod{12}$.
\end{solution}

\begin{solution}{ex:primcount}
\idea Sort the characters modulo $q$ by their conductor. Each conductor $d$ contributes exactly its primitive characters, so $\varphi$ is the
divisor sum of $\varphi^*$.
\writeup By Proposition~\ref{prop:conductor}(c), $\chi\mapsto(d^*,\chi^*)$ maps the characters modulo $q$ to pairs ($d\mid q$, primitive character modulo $d$). It is
injective, since $\chi^*$ determines $\chi$ (they agree on integers prime to $q$, and $\chi$ vanishes elsewhere). It is surjective: a primitive $\psi$
modulo $d\mid q$ induces the character $\chi=\chi_0\psi$ modulo $q$, and by uniqueness in (c), $\chi^*=\psi$. Counting gives $\varphi(q)=\sum_{d\mid q}\varphi^*(d)$, so
$\varphi^*=\mu*\varphi$ by Möbius inversion (Tome~1, Theorem~4.4), which is multiplicative as a convolution of multiplicative functions. At $p^k$:
$\varphi^*(p)=\varphi(p)-\varphi(1)=p-2$ and, for $k\ge2$, $\varphi^*(p^k)=\varphi(p^k)-\varphi(p^{k-1})=p^{k-2}(p-1)^2=p^k(1-\frac1p)^2$. In particular $\varphi^*(2)=0$, so if
$q=2m$ with $m$ odd, $\varphi^*(q)=\varphi^*(2)\varphi^*(m)=0$. Modulo $12$ this gives $\varphi^*(4)\varphi^*(3)=1\cdot1=1$, matching Exercise~\ref{ex:chartables}.
\end{solution}

\begin{solution}{ex:gausssquare}
\idea Conjugation inverts $e(\frac aq)$ and $\chi(a)$; reflecting $a\mapsto-a$ turns the conjugate sum back into a Gauss sum, at the cost of $\chi(-1)$.
\writeup Substituting $b=-a$,
\[
\begin{aligned}\overline{\tau(\chi)}&=\sum_a\overline\chi(a)e\Bigl(-\frac aq\Bigr)=\sum_b\overline\chi(-b)e\Bigl(\frac bq\Bigr)\eqbr=\overline\chi(-1)\tau(\overline\chi),\end{aligned}
\]
and $\overline\chi(-1)=\chi(-1)$ since $\chi(-1)=\pm1$. So, by Theorem~\ref{thm:gauss},
\[
\begin{aligned}\tau(\chi)\tau(\overline\chi)&=\chi(-1)\tau(\chi)\overline{\tau(\chi)}\eqbr=\chi(-1)|\tau(\chi)|^2=\chi(-1)q.\end{aligned}
\]
For real $\chi$, $\overline\chi=\chi$. Check: $\tau(\chi_{-4})^2=(2i)^2=-4$.
\end{solution}

\begin{solution}{ex:lpositive}
\idea On the real axis $L(\sigma,\chi)$ is real and positive for $\sigma>1$; it cannot become negative at $\sigma=1$ without passing through $0$, which
Theorem~\ref{thm:lnonzero} forbids.
\writeup For $\sigma>1$, $L(\sigma,\chi)=\exp(S)$ with $S$ real (the coefficients $\frac{\chi(p)^k}{kp^{k\sigma}}$ are real), so $L(\sigma,\chi)>0$. The function $\sigma\mapsto L(\sigma,\chi)$ is
continuous on $\left]0,\infty\right[$ (Theorem~\ref{thm:lbasic}(b)), so $L(1,\chi)=\lim_{\sigma\to1^+}L(\sigma,\chi)\ge0$, and $L(1,\chi)\ne0$.
\end{solution}

\begin{solution}{ex:l1values}
\idea Both characters are real and primitive, so the work is to compute the Gauss sum and a short finite sum.
\writeup $\chi_{-3}$ is odd, primitive, with $\tau(\chi_{-3})=e(\frac13)-e(\frac23)=2i\sin\frac{2\pi}3=i\sqrt3$. The first formula gives
$L(1,\chi_{-3})=-\frac{i\pi}{3\cdot i\sqrt3}(1\cdot1-1\cdot2)=\frac\pi{3\sqrt3}=0.6046\ldots$. $\chi_5$ is even, primitive, with $\tau(\chi_5)=2\cos\frac{2\pi}5-2\cos\frac{4\pi}5=\frac{\sqrt5-1}2+\frac{\sqrt5+1}2=\sqrt5$
(pairing $a$ with $5-a$). Since $\sin\frac{4\pi}5=\sin\frac\pi5$ and $\sin\frac{3\pi}5=\sin\frac{2\pi}5$, the second formula gives
\[
\begin{aligned}L(1,\chi_5)&=-\frac1{\sqrt5}\Bigl(2\log\sin\frac\pi5-2\log\sin\frac{2\pi}5\Bigr)\eqbr=\frac2{\sqrt5}\log\frac{\sin(2\pi/5)}{\sin(\pi/5)}.\end{aligned}
\]
And $\frac{\sin(2\pi/5)}{\sin(\pi/5)}=2\cos\frac\pi5=\frac{1+\sqrt5}2$, the golden ratio, so $L(1,\chi_5)=\frac2{\sqrt5}\log\frac{1+\sqrt5}2=0.4304\ldots$; the notebook checks both against the
series. The golden ratio is the fundamental unit of $\Z[\frac{1+\sqrt5}2]$: the even case of Dirichlet's class number formula.
\end{solution}

\begin{solution}{ex:residuesum}
\idea For odd real $\chi$, Dirichlet's formula expresses the positive number $L(1,\chi)$ as a negative multiple of $\sum\chi(a)a$, once Gauss's sign is
known. So $\sum\chi(a)a<0$, which is the claim.
\writeup Let $\chi=\bigl(\frac\cdot p\bigr)$, odd since $p\equiv3\pmod4$, and primitive. By Tome~4, Theorem~8.14, $\tau(\chi)=i\sqrt p$. The first formula of Theorem~\ref{thm:l1formula}
gives $L(1,\chi)=-\frac{i\pi}{p\cdot i\sqrt p}\sum_a\chi(a)a=-\frac\pi{p^{3/2}}\sum_a\chi(a)a$. By Exercise~\ref{ex:lpositive}, $L(1,\chi)>0$, so $\sum_a\chi(a)a<0$: the non-residues add up to more
than the residues. For $p=7$ the residues are $1,2,4$ (sum $7$) and the non-residues $3,5,6$ (sum $14$); for $p=11$, residues $1,3,4,5,9$ (sum $22$),
non-residues $2,6,7,8,10$ (sum $33$). The difference divided by $p$ is $1$ in both cases: $h(-7)=h(-11)=1$. No proof avoiding analysis is known.
\end{solution}

\begin{solution}{ex:lderivative}
\idea Near $s=1$ the first $q$ terms of the series carry the size, at most a harmonic sum; the rest is controlled by the bounded partial sums
through Abel summation.
\writeup By the proof of Tome~4, Proposition~10.1(b), the partial sums of $L(s,\chi)$ converge uniformly on closed discs in $\sigma>0$, so by Tome~4,
Theorem~9.15(a) the series may be differentiated term by term. Let $A(t)=\sum_{n\le t}\chi(n)$, so $|A(t)|\le\varphi(q)\le q$ and $A(q)=0$. For a $C^1$ function $f$
decreasing to $0$ on $[q,\infty[$, Lemma~\ref{lem:abelint} gives $\sum_{q<n\le x}\chi(n)f(n)=A(x)f(x)-\int_q^xA(t)f'(t)\dd t$, so the tail $\sum_{n>q}\chi(n)f(n)$ has modulus at most
$\varphi(q)\int_q^\infty(-f')=\varphi(q)f(q)$.

With $f(t)=\frac1t$: the head is $\bigl|\sum_{n<q}\frac{\chi(n)}n\bigr|\le1+\log q$ (Tome~4, Corollary~11.6, or comparison with $\int_1^q\frac{\mathrm dt}t$), and the tail is at most
$\frac{\varphi(q)}q\le1$. So $|L(1,\chi)|\le\log q+2$.

With $f(t)=\frac{\log t}t$, decreasing on $[e,\infty[\,\supseteq[q,\infty[$: the tail is at most $\varphi(q)\frac{\log q}q\le\log q$. For the head, $\frac{\log n}n\le\int_{n-1}^n\frac{\log t}t\dd t$ for
$n\ge4$, so $\sum_{n<q}\frac{\log n}n\le\frac{\log2}2+\frac{\log3}3+\int_3^q\frac{\log t}t\dd t\le1+\frac12\log^2q$. Adding up gives the bound for $L'(1,\chi)$. These bounds are crude
but uniform in $\chi$; Chapter~4 needs only that $L(1,\chi)$ and $L'(1,\chi)$ are finite, and Chapter~11 sharpens such estimates.
\end{solution}

\section*{Sources and further reading}

\begin{itemize}
\item P.~G.~L.~Dirichlet, \emph{Beweis des Satzes, dass jede unbegrenzte arithmetische Progression, deren erstes Glied und Differenz ganze
Zahlen ohne gemeinschaftlichen Factor sind, unendlich viele Primzahlen enthält}, Abhandlungen der Königlichen Preussischen Akademie
der Wissenschaften (1837), 45--81.
\item H.~Davenport, \emph{Multiplicative Number Theory}, 3rd ed., revised by H.~L.~Montgomery, Springer GTM~74, 2000, Chapters~1, 2 and~4--6
(characters, Gauss sums, $L(1,\chi)\ne0$, primitive characters, the class number formula).
\item T.~M.~Apostol, \emph{Introduction to Analytic Number Theory}, Springer, 1976, Chapters~6--8 (characters, Dirichlet's theorem, Gauss sums, induced
characters).
\item J.-P.~Serre, \emph{A Course in Arithmetic}, Springer GTM~7, 1973, Chapter~VI (Dirichlet's theorem with Landau's lemma on series with
positive coefficients, the route followed here).
\item J.~Elstrodt, \emph{The life and work of Gustav Lejeune Dirichlet (1805--1859)}, Clay Mathematics Proceedings~7 (2007), 1--37. The source
of the epigraph.
\end{itemize}
